\documentclass[11pt]{article}
\usepackage[margin=1in]{geometry}
\usepackage{amsmath,amssymb,amsthm,mathtools,bm}
\usepackage{microtype}
\usepackage{hyperref}
\usepackage{enumitem}
\usepackage{booktabs}
\usepackage{xcolor}
\hypersetup{colorlinks=true,linkcolor=blue,citecolor=blue,urlcolor=blue}

\newcommand{\E}{\mathbb{E}}
\newcommand{\R}{\mathbb{R}}
\newcommand{\Law}{\mathrm{Law}}
\newcommand{\divergence}{\nabla\!\cdot\!}
\newcommand{\tr}{\mathrm{tr}}
\newcommand{\sg}{\mathrm{sg}}
\newcommand{\dd}{\mathrm{d}}
\newcommand{\I}{\mathrm{I}}
\newcommand{\norm}[1]{\left\lVert #1 \right\rVert}
\newcommand{\inner}[2]{\left\langle #1, #2 \right\rangle}
\newcommand{\abs}[1]{\left| #1 \right|}
\newcommand{\grad}{\nabla}
\newcommand{\J}{\mathbf{J}}
\newcommand{\flow}{X_{\psi}}
\newcommand{\energy}{U_{\phi}}

\title{A Two-Head Recipe for Learning a Cross-Mode Consistent Energy\vspace{0.2em}\\
\large Invertible few-step flow head + energy head + occasional exact likelihood anchoring}
\author{}
\date{}

\begin{document}
\maketitle

\begin{abstract}
The core idea is to let an invertible few-step flow head provide \emph{absolute} likelihood information, and let an energy head learn the same family of marginals through score matching. The flow head fixes the cross-mode additive constants that ordinary score matching cannot see, while the energy head supplies a conservative score field and a cheap scalar energy at test time. The cleanest formulation is to define both heads on the same stochastic-interpolant time variable. Then the few-step head is trained with a Lagrangian map loss plus invertibility and semigroup regularization, and the energy head is trained with score matching plus an occasional exact value-matching loss to the flow-induced likelihood. A useful additional cross-head regularizer comes from the Liouville equation: if the flow head induces velocity $v_t$, then the normalized energy must satisfy $\partial_t U + \inner{v_t}{\grad U} = \divergence v_t$. This gives a local MSE constraint that transports normalization across modes.
\end{abstract}

\section{Why this should solve the mode-consistency problem}

Let $p_{\mathrm{data}}$ be the data distribution on $\R^d$, and let $p_0$ be a tractable base density, typically $\mathcal{N}(0,I)$. The failure mode of ordinary score matching is that if the data are multimodal, one may learn a very good score field on each mode while still missing the \emph{relative additive constants} between modes. In other words, one learns energies only up to mode-dependent offsets. The whole point of the proposed hybrid model is that the flow head has access to a normalized density through change of variables, so it can transmit those missing offsets to the energy head.

The design principle is therefore:
\begin{enumerate}[label=(\roman*),itemsep=0.25em,topsep=0.25em]
    \item train a two-time invertible map $\flow(x,s,t)$ that approximates the probability flow between times $s$ and $t$;
    \item use this map to compute an absolute log density by pulling a point back to time $0$ and correcting by a Jacobian determinant;
    \item train a scalar energy $\energy(x,t)$ whose gradient is a score, but occasionally anchor its value to the exact likelihood supplied by the map head.
\end{enumerate}

That is the mechanism that should remove cross-mode inconsistency.

\section{Joint stochastic-interpolant model}

I would put both heads on the same time-indexed family of marginals. Let
\begin{equation}
I_t = \alpha_t x_0 + \beta_t x_1 + \gamma_t \varepsilon,
\qquad x_0 \sim p_0,
\quad x_1 \sim p_{\mathrm{data}},
\quad \varepsilon \sim \mathcal{N}(0,I),
\label{eq:interpolant}
\end{equation}
with $t \in [0,1]$, $\alpha_0=1$, $\beta_1=1$, $\alpha_1=\beta_0=0$, and $\gamma_0=\gamma_1=0$. Let $\rho_t = \Law(I_t)$.

We learn two heads:
\begin{itemize}[itemsep=0.25em,topsep=0.25em]
    \item a \textbf{few-step flow-map head} $\flow(\cdot,s,t): \R^d \to \R^d$, parameterized as
    \begin{equation}
    \flow(x,s,t) = x + (t-s) h_{\psi}(x,s,t),
    \label{eq:flow_parameterization}
    \end{equation}
    so that $\flow(x,t,t)=x$ by construction;
    \item an \textbf{energy head} $\energy(x,t) \approx -\log \rho_t(x)$ whose score is
    \begin{equation}
    s_{\phi}(x,t) := \grad_x \energy(x,t).
    \label{eq:score_from_energy}
    \end{equation}
\end{itemize}

At $t=1$, the target is the data energy $-\log p_{\mathrm{data}}(x)$; at intermediate $t$, the target is the smoothed/interpolated marginal energy $-\log \rho_t(x)$.

\section{Few-step head: internal consistency losses}

\subsection{Lagrangian loss (recommended primary map loss)}

Among the available consistency objectives, the Lagrangian one is the cleanest and usually the one I would trust most. A teacher-free direct objective is obtained by requiring the map to reproduce the time derivative of the interpolant after pulling $I_t$ back to time $s$:
\begin{equation}
\mathcal{L}_{\mathrm{Lag}}
=
\E_{s,t,x_0,x_1,\varepsilon}
\Big[
\big\| \partial_t \flow\big(\flow(I_t,t,s),s,t\big) - \dot I_t \big\|^2
\Big].
\label{eq:L_lag}
\end{equation}
This is the direct map-matching analogue of the Lagrangian consistency objective.

\subsection{Invertibility loss}

To make likelihood evaluation possible, the map should be explicitly encouraged to be invertible before convergence. The practical loss is the cycle loss
\begin{equation}
\mathcal{L}_{\mathrm{inv}}
=
\E_{s,t,x}
\Big[
\big\| x - \flow(\flow(x,s,t),t,s) \big\|^2
\Big].
\label{eq:L_inv}
\end{equation}
This is much better than a generic ``Jacobian penalty'' as the primary mechanism, because it directly enforces the existence of an inverse map. If one wants an additional Jacobian-conditioning term, it should be light and auxiliary, not the main invertibility device.

\subsection{Semigroup loss}

The two-time map should also satisfy composition:
\begin{equation}
\flow\big(\flow(x,r,s),s,t\big) = \flow(x,r,t).
\label{eq:semigroup_identity}
\end{equation}
So an auxiliary semigroup loss is
\begin{equation}
\mathcal{L}_{\mathrm{sg}}
=
\E_{r,s,t,x}
\Big[
\big\| \flow(\flow(x,r,s),s,t) - \flow(x,r,t) \big\|^2
\Big].
\label{eq:L_sg}
\end{equation}
This is especially useful if one wants the same trained head to work with different inference grids.

\subsection{Eulerian loss (interesting, but I would downweight it)}

An Eulerian formulation penalizes the PDE residual at a fixed spatial point. In direct form one may write
\begin{equation}
\mathcal{L}_{\mathrm{Eul}}
=
\E_{s,t,x_0,x_1,\varepsilon}
\Big[
\big\| \partial_s \flow(I_s,s,t) - \sg\big( \dot I_s \cdot \grad_x \flow(I_s,s,t) \big) \big\|^2
\Big].
\label{eq:L_eul}
\end{equation}
This is mathematically interesting, but it depends on spatial Jacobians of the map and is usually less attractive numerically. My recommendation is:
\begin{equation}
\text{main map loss } = \mathcal{L}_{\mathrm{Lag}} + \lambda_{\mathrm{inv}} \mathcal{L}_{\mathrm{inv}} + \lambda_{\mathrm{sg}} \mathcal{L}_{\mathrm{sg}} + \lambda_{\mathrm{Eul}} \mathcal{L}_{\mathrm{Eul}},
\qquad \lambda_{\mathrm{Eul}} \ll 1.
\label{eq:map_total}
\end{equation}

\section{Exact likelihood coming from the flow head}

Once the flow head is approximately invertible, it defines a normalized density. Let
\begin{equation}
B_{\psi,t}(x) := \flow(x,t,0)
\label{eq:backward_map}
\end{equation}
be the backward map to the base. Then the induced flow-head density is
\begin{equation}
q_{\psi,t}(x) = p_0\big(B_{\psi,t}(x)\big)
\,\abs{\det \J_x B_{\psi,t}(x)}.
\label{eq:flow_density_inverse}
\end{equation}
Equivalently, if $x_t = \flow(z,0,t)$ with $z \sim p_0$, then
\begin{equation}
q_{\psi,t}(x_t)
=
\frac{p_0(z)}{\abs{\det \J_z \flow(z,0,t)}}.
\label{eq:flow_density_forward}
\end{equation}
So the flow head supplies an \emph{absolute} teacher energy
\begin{equation}
\widetilde U_{\psi}(x,t)
:= -\log q_{\psi,t}(x)
= -\log p_0\big(B_{\psi,t}(x)\big) - \log \abs{\det \J_x B_{\psi,t}(x)}.
\label{eq:teacher_energy_exact}
\end{equation}
This is the critical object: it contains the missing cross-mode constants.

\section{Energy head: local score losses}

\subsection{Space score matching on the interpolant family}

Write
\begin{equation}
\mu_t := \alpha_t x_0 + \beta_t x_1,
\qquad I_t = \mu_t + \gamma_t \varepsilon.
\label{eq:mu_t}
\end{equation}
Conditioned on $(x_0,x_1)$, the density of $I_t$ is Gaussian with mean $\mu_t$ and covariance $\gamma_t^2 I$. Therefore
\begin{equation}
\grad_x \log \rho_t(x)
=
\E\left[
-\frac{x-\mu_t}{\gamma_t^2}
\;\middle|\; I_t=x
\right],
\label{eq:score_identity}
\end{equation}
which implies the negative-score identity
\begin{equation}
\grad_x U^{\star}(x,t)
=
\E\left[
\frac{x-\mu_t}{\gamma_t^2}
\;\middle|\; I_t=x
\right],
\qquad U^{\star}(x,t)=-\log \rho_t(x).
\label{eq:negative_score_identity}
\end{equation}
Hence a direct MSE training loss for the energy head is
\begin{equation}
\mathcal{L}_{\mathrm{SM}}
=
\E_{t,x_0,x_1,\varepsilon}
\left[
 w_{\mathrm{SM}}(t)
\left\|
\grad_x \energy(I_t,t) - \frac{I_t-\mu_t}{\gamma_t^2}
\right\|^2
\right].
\label{eq:L_SM}
\end{equation}
A natural scale-balancing choice is $w_{\mathrm{SM}}(t)=\gamma_t^2/d$.

\subsection{Optional time-score matching on the interpolant family}

There is also a direct time-derivative identity. Since
\begin{equation}
\log p(I_t\mid x_0,x_1)
=
-\frac{d}{2}\log(2\pi\gamma_t^2)
-
\frac{\norm{I_t-\mu_t}^2}{2\gamma_t^2},
\label{eq:conditional_log_gauss}
\end{equation}
we get
\begin{equation}
\partial_t \log p(x\mid x_0,x_1)
=
-d\frac{\dot\gamma_t}{\gamma_t}
+
\frac{\dot\gamma_t}{\gamma_t^3}\norm{x-\mu_t}^2
+
\frac{\inner{x-\mu_t}{\dot\mu_t}}{\gamma_t^2},
\label{eq:conditional_time_derivative}
\end{equation}
where $\dot\mu_t = \dot\alpha_t x_0 + \dot\beta_t x_1$. Therefore
\begin{equation}
\partial_t U^{\star}(x,t)
=
\E\left[
 d\frac{\dot\gamma_t}{\gamma_t}
-
\frac{\dot\gamma_t}{\gamma_t^3}\norm{x-\mu_t}^2
-
\frac{\inner{x-\mu_t}{\dot\mu_t}}{\gamma_t^2}
\;\middle|\; I_t=x
\right].
\label{eq:time_identity}
\end{equation}
This gives an optional transported analogue of dual score matching:
\begin{equation}
\mathcal{L}_{\mathrm{TM}}
=
\E_{t,x_0,x_1,\varepsilon}
\left[
 w_{\mathrm{TM}}(t)
\left(
\partial_t \energy(I_t,t)
-
\left[
 d\frac{\dot\gamma_t}{\gamma_t}
-
\frac{\dot\gamma_t}{\gamma_t^3}\norm{I_t-\mu_t}^2
-
\frac{\inner{I_t-\mu_t}{\dot\mu_t}}{\gamma_t^2}
\right]
\right)^2
\right].
\label{eq:L_TM}
\end{equation}
I would treat this as optional. The exact cross-head likelihood anchoring below is stronger.

\section{Cross-head consistency: the key innovation}

\subsection{Occasional exact value matching (main cross-head loss)}

The simplest and strongest way to let the energy head ``see'' normalized likelihood information is just to regress to the exact scalar energy supplied by the flow head:
\begin{equation}
\mathcal{L}_{\mathrm{val}}
=
\E_{t,x \sim \nu_t}
\Big[
 w_{\mathrm{val}}(t)
\big(
\energy(x,t) - \sg(\widetilde U_{\psi}(x,t))
\big)^2
\Big],
\label{eq:L_val}
\end{equation}
where $\nu_t$ may be either the flow-head marginal $q_{\psi,t}$, the interpolant marginal $\rho_t$, or a mixture of the two. The stop-gradient is important: during this step, the likelihood should serve as a teacher for the energy head, not as a signal that destabilizes the map head.

This is the loss that directly fixes cross-mode additive constants. It should \emph{not} be run on every batch. It is expensive because it needs Jacobian determinants (or log-determinants accumulated over a discrete path), so I would turn it on only every $K$ optimization steps, or with Bernoulli probability $p_{\mathrm{anchor}} \in [0.01,0.1]$.

\subsection{Liouville transport loss (my favorite auxiliary cross-head loss)}

There is also a very clean PDE identity. Let $v_{\psi}(x,t)$ denote the instantaneous velocity induced by the flow head, for example
\begin{equation}
 v_{\psi}(x,t) := \partial_{\tau}\flow(x,t,\tau)\big|_{\tau=t}.
\label{eq:velocity_from_map}
\end{equation}
If $q_{\psi,t}$ is the density induced by the flow head, then it satisfies the continuity equation
\begin{equation}
\partial_t q_{\psi,t} + \divergence(q_{\psi,t} v_{\psi}) = 0.
\label{eq:continuity}
\end{equation}
Setting $\widetilde U_{\psi}(x,t) = -\log q_{\psi,t}(x)$ and expanding \eqref{eq:continuity} gives the transport identity
\begin{equation}
\partial_t \widetilde U_{\psi}(x,t)
+
\inner{v_{\psi}(x,t)}{\grad_x \widetilde U_{\psi}(x,t)}
=
\divergence v_{\psi}(x,t).
\label{eq:liouville_teacher}
\end{equation}
So a local cross-head residual is
\begin{equation}
\mathcal{L}_{\mathrm{Liou}}
=
\E_{t,x \sim \nu_t}
\Big[
 w_{\mathrm{Liou}}(t)
\big(
\partial_t \energy(x,t)
+
\inner{v_{\psi}(x,t)}{\grad_x \energy(x,t)}
-
\divergence v_{\psi}(x,t)
\big)^2
\Big].
\label{eq:L_Liou}
\end{equation}
This is a beautiful constraint because it transports normalization information \emph{locally in time} and is the exact multivariable-calculus bridge between the two heads. It says: along a characteristic of the map head, the energy must increase at rate $\divergence v_{\psi}$.

In practice, $\divergence v_{\psi}$ can be estimated with a Hutchinson trace estimator if $d$ is large. This makes \eqref{eq:L_Liou} far cheaper than computing exact log-determinants on every step.

\subsection{Boundary calibration at the base}

If one uses the Liouville residual, one should also pin the boundary condition at $t=0$:
\begin{equation}
\mathcal{L}_{\mathrm{base}}
=
\E_{z \sim p_0}
\Big[
\big(
\energy(z,0) + \log p_0(z)
\big)^2
\Big].
\label{eq:L_base}
\end{equation}
Together, \eqref{eq:L_Liou} and \eqref{eq:L_base} say that the energy head should solve the same transport PDE as the normalized flow-head likelihood, starting from the known base energy.

\section{The total loss I would actually use}

A practical combined objective is
\begin{align}
\mathcal{L}_{\mathrm{flow}}
&=
\mathcal{L}_{\mathrm{Lag}}
+
\lambda_{\mathrm{inv}}\mathcal{L}_{\mathrm{inv}}
+
\lambda_{\mathrm{sg}}\mathcal{L}_{\mathrm{sg}}
+
\lambda_{\mathrm{Eul}}\mathcal{L}_{\mathrm{Eul}},
\label{eq:L_flow_total}
\\
\mathcal{L}_{\mathrm{energy}}
&=
\lambda_{\mathrm{SM}}\mathcal{L}_{\mathrm{SM}}
+
\lambda_{\mathrm{TM}}\mathcal{L}_{\mathrm{TM}}
+
\lambda_{\mathrm{base}}\mathcal{L}_{\mathrm{base}}
+
\lambda_{\mathrm{Liou}}\mathcal{L}_{\mathrm{Liou}}
+
\mathbf{1}_{\mathrm{anchor\ batch}}\,\lambda_{\mathrm{val}}\mathcal{L}_{\mathrm{val}},
\label{eq:L_energy_total}
\\
\mathcal{L}_{\mathrm{joint}}
&=
\mathcal{L}_{\mathrm{flow}} + \mathcal{L}_{\mathrm{energy}}.
\label{eq:L_joint}
\end{align}

My recommended default is:
\begin{itemize}[itemsep=0.25em,topsep=0.25em]
    \item make $\mathcal{L}_{\mathrm{Lag}}$ the main flow-head loss;
    \item always keep $\mathcal{L}_{\mathrm{inv}}$ on;
    \item use $\mathcal{L}_{\mathrm{sg}}$ at moderate weight if you want robust multistep closure;
    \item keep $\mathcal{L}_{\mathrm{Eul}}$ tiny or off;
    \item train the energy head mostly with $\mathcal{L}_{\mathrm{SM}} + \mathcal{L}_{\mathrm{Liou}} + \mathcal{L}_{\mathrm{base}}$;
    \item turn on $\mathcal{L}_{\mathrm{val}}$ only occasionally to recalibrate the absolute energy scale and cross-mode offsets.
\end{itemize}

\section{Training schedule}

The most stable schedule is staged.

\paragraph{Stage 1: map-head warmup.}
Train the flow head with \eqref{eq:L_flow_total} until the inverse-cycle error and semigroup error are small enough that likelihoods are numerically meaningful.

\paragraph{Stage 2: energy-head warmup.}
Train the energy head with score matching, i.e. \eqref{eq:L_SM}, optionally \eqref{eq:L_TM}, plus the base condition \eqref{eq:L_base}.

\paragraph{Stage 3: joint calibration.}
Turn on the cross-head terms. Most steps use the cheap local transport residual \eqref{eq:L_Liou}. Every $K$ steps, run an anchor batch with the exact scalar target \eqref{eq:L_val}. That is where the energy head receives the absolute likelihood information that ordinary score matching cannot provide.

\section{Why the proposal is internally coherent}

The flow head and energy head are not learning unrelated objects.

\begin{enumerate}[label=(\arabic*),itemsep=0.25em,topsep=0.25em]
    \item The flow head defines a normalized density $q_{\psi,t}$ through change of variables.
    \item The energy head is trained so that its gradient matches the score of the same marginal family.
    \item The Liouville residual forces the scalar energy to evolve in time exactly as a log density must evolve under the flow-head velocity field.
    \item The occasional value anchor identifies the global additive constant and transfers the correct relative offsets between disconnected or weakly connected modes.
\end{enumerate}

This is precisely the ingredient missing from pure score matching.

\section{What I would and would not do}

\paragraph{What I would do.}
I would use \eqref{eq:L_val} as the central innovation, because it is the direct supervision signal that makes the energy normalized and cross-mode consistent. I would support it with \eqref{eq:L_Liou}, because it is a principled and cheap local proxy for the same information. I would use the Lagrangian map loss as the main internal consistency loss, with semigroup as the main auxiliary.

\paragraph{What I would avoid.}
I would not rely on a bare Jacobian penalty alone to enforce invertibility. I would also avoid making the Eulerian loss a major term: it is elegant, but it pushes spatial Jacobians of the map directly into the objective and is usually the least attractive numerically. Finally, I would not run the exact likelihood anchor on every batch: it is too expensive and unnecessary once the model is roughly calibrated.

\section{A short theorem-like summary}

Suppose the flow head is invertible and induces a smooth velocity field $v_{\psi}$. Then the normalized teacher energy $\widetilde U_{\psi}=-\log q_{\psi,t}$ is characterized by the pair
\begin{equation}
\widetilde U_{\psi}(x,0) = -\log p_0(x),
\qquad
\partial_t \widetilde U_{\psi} + \inner{v_{\psi}}{\grad \widetilde U_{\psi}} = \divergence v_{\psi}.
\label{eq:teacher_PDE_summary}
\end{equation}
Therefore, if the energy head achieves small score-matching error and small values of $\mathcal{L}_{\mathrm{base}}$, $\mathcal{L}_{\mathrm{Liou}}$, and occasional $\mathcal{L}_{\mathrm{val}}$, then it is being forced toward the same normalized energy that is implied by the invertible flow head. That is the mathematical reason the construction should remove mode-dependent additive constants.

\section{Bottom line}

Yes: the idea is sound, and there is a very natural way to formalize it.

The best version is not just ``flow head + scalar MSE.'' It is:
\begin{enumerate}[label=(\arabic*),itemsep=0.25em,topsep=0.25em]
    \item a two-time invertible few-step map head trained mainly by a Lagrangian consistency loss plus invertibility and semigroup regularization;
    \item an energy head trained by score matching on the same interpolant family;
    \item an occasional exact likelihood distillation loss from the flow head to the energy head;
    \item a local Liouville transport loss that ties the time evolution of the energy to the velocity and divergence of the flow head.
\end{enumerate}

That gives you exactly what you want: a conservative score field with the global cross-mode constants supplied by a normalized few-step model.

\end{document}
